\section{Conclusion}
\label{sec:conclusion}

This paper presented \texttt{sionnart}: a pybind11-embedded integration of the Sionna~RT ray tracing engine into ns-3, augmented by a displacement-based coherence cache and a proactive adaptive virtual receiver mechanism, both designed to keep per-step GPU overhead manageable.
The embedded design achieves up to $232\times$ lower runtime than the two-process \texttt{ns3sionna} baseline, maintains a GPU memory footprint below 500~MiB, and in stationary scenarios outperforms even the analytical Friis baseline owing to cache-driven elimination of repeated ray tracing calls.
Applied to a 3GPP~NR deployment across free-space, urban macro, and urban micro environments, the framework shows that a geometry-aware channel model produces more conservative — and physically more detailed — channel quality estimates than stochastic models, and confirms that $8\times8$ gNB antenna arrays are essential for overcoming harsh propagation conditions.
As future work, the site-specific multipath data produced by \texttt{sionnart} will be used to extend the framework toward Integrated Sensing and Communications (ISAC) evaluation, enabling joint simulation of communication and sensing within a single ns-3 experiment.